%! Observational Studies — Unit 1, Lesson 1.5 — Slides (Beamer)
% PILOT SHAPE (2026-09-06): modeled on AP Statistics lesson 1.4 minus its AP Practice page.
% GRADUAL RELEASE deck: title → targets → warm-up → hook (dark, unresolved) → I DO divider
% (dark) → two or three frames per notes ROW, labelled by row name and problem numbers (the
% crux frame flagged \sectionlabel[redacc]{}) → Guided Practice (a live surface, un-answered)
% → spoken debrief with the cold check → close & START THE HOMEWORK, which is the individual
% practice → close (dark). No group activity, no practice launch, no exit ticket. There is no
% spoiler rule: the notes frames ARE the direct instruction.
% \CourseName is NOT defined in beamer — the course name is written literally.
% saar-beamer does not load tcolorbox (use block) and beamer's itemize takes no [leftmargin=].
\documentclass[aspectratio=169, 11pt]{beamer}
\usepackage{saar-beamer}   % provides \ceruleanheader{} and \sectionlabel[color]{}

\newcommand{\redink}[1]{{\color{redacc}\bfseries #1}}

\begin{document}

% TITLE SLIDE (CourseName is not defined in beamer, so the name is literal)
{
\setbeamercolor{background canvas}{bg=cerulean}
\begin{frame}[plain]
  \vspace{2.8cm}\hspace{0.55cm}%
  \begin{minipage}{0.9\textwidth}
    {\color{goldacc}\bfseries\small LESSON 1.5}\\[0.5cm]
    {\color{white}\bfseries\LARGE Observational Studies}\\[1.2cm]
    {\color{frostmid}\small Statistical Analysis \& Algebraic Reasoning~$\cdot$~Unit 1}
  \end{minipage}
\end{frame}
}

% ── LEARNING TARGETS ─────────────────────────────────────────────────────────
\begin{frame}[t, plain]
  \ceruleanheader{Today's targets}
  \sectionlabel{Lesson 1.5}
  By the end of today, I can\ldots
  \begin{itemize}\setlength{\itemsep}{5pt}
    \item explain what makes a study \textbf{observational} --- nobody is \emph{assigned} to a
          group --- and plan one with the \textbf{statistical cycle}.
    \item compare two groups that already exist and name the \textbf{explanatory} and
          \textbf{response} variables.
    \item describe the result as an \textbf{association}, with no causal verb.
    \item explain why an association is not \textbf{causation} --- naming a reversed arrow or a
          \textbf{lurking variable}.
  \end{itemize}
  \begin{block}{How today works}
    Warm-up $\to$ \textbf{notes together, row by row} $\to$ \textbf{one study we work as a
    class} $\to$ debrief $\to$ \textbf{you start the homework here, alone}.
  \end{block}
\end{frame}

% ── WARM-UP ──────────────────────────────────────────────────────────────────
\begin{frame}[t, plain]
  \ceruleanheader{Warm-Up --- 5 minutes, silent}
  \sectionlabel{back to the Harbor Point Community Pool}
  \vspace{0.3cm}
  \begin{center}
    \begin{minipage}{0.88\textwidth}
      A random sample of \textbf{200} of the pool's \textbf{1{,}500} members --- \textbf{80}
      swim in the morning, \textbf{120} in the evening. Everyone answered.
      \begin{enumerate}\setlength{\itemsep}{7pt}
        \item \textbf{56} of the \textbf{80} morning swimmers get 7+ hours of sleep. Percent?
        \item \textbf{54} of the \textbf{120} evening swimmers do. Percent?
        \item The gap in percentage points --- then put the whole sample back together.
      \end{enumerate}
    \end{minipage}
  \end{center}
  \vspace{0.35cm}
  \begin{center}
    {\color{cerulean}\bfseries Two groups, two different wholes. 70\%, 45\%, 25 points, 55\%
    --- these stay on the board all period.}
  \end{center}
\end{frame}

% ── HOOK — leave it UNRESOLVED; problem 13 (row 4, Not Causation) settles it ─
{
\setbeamercolor{background canvas}{bg=cerulean}
\begin{frame}[plain]
  \vspace{1.1cm}
  \begin{center}
    {\color{frostmid}\bfseries\small RANDOM SAMPLE OF 200. EVERYONE ANSWERED. EVERY COUNT CORRECT.}\\[0.7cm]
    {\color{white}\bfseries\Large 70\% of morning swimmers get 7+ hours of sleep.\\[0.25cm]
    Only 45\% of evening swimmers do.}\\[0.7cm]
    {\color{goldacc}\bfseries\large The flyer: ``Swim in the morning --- sleep better!''}\\[0.4cm]
    {\color{white}\bfseries\large Is the claim safe?}\\[0.35cm]
    {\color{frostmid}\small Circle \emph{safe} or \emph{not safe} in your hook box. We settle
    this at problem 13.}
  \end{center}
\end{frame}
}

% ── I DO — direct instruction begins ─────────────────────────────────────────
{
\setbeamercolor{background canvas}{bg=cerulean}
\begin{frame}[plain]
  \vspace{1.8cm}\hspace{0.8cm}%
  \begin{minipage}{0.86\textwidth}
    {\color{goldacc}\bfseries\small GUIDED NOTES \ \textbullet\ \ I DO}\\[0.7cm]
    {\color{white}\bfseries\Large Open to your Guided Notes page. We work the table row by
    row --- fill each term as we name it, not before.}\\[0.9cm]
    {\color{frostmid}\small Five terms go in the vocabulary box today. Write each one the
    moment we use it.}
  \end{minipage}
\end{frame}
}

% ── ROW 1a. ONE WORD: ASSIGN ─────────────────────────────────────────────────
\begin{frame}[t, plain]
  \ceruleanheader{One word: assign}
  \sectionlabel{notes row 1 $\cdot$ Assign $\cdot$ problems 1--4}
  \vspace{4pt}
  \begin{columns}[t]
    \begin{column}{0.47\textwidth}
      \begin{block}{Watch --- observational}
        The members already swim when they like. You ask 200 of them when they swim and how
        they sleep. \textbf{The members decided.}
      \end{block}
    \end{column}
    \begin{column}{0.47\textwidth}
      \begin{block}{Assign --- an experiment (1.6)}
        The director \emph{puts} 100 members at 6 a.m.\ and 100 at 7 p.m.\ for a month.
        \textbf{The director decided.}
      \end{block}
    \end{column}
  \end{columns}
  \vspace{10pt}
  \begin{center}
    {\large In an \textbf{observational study} you measure or survey \redink{without assigning}
    anybody to a group.}\\[6pt]
    {\small\color{slate} \textbf{A survey is one type of observational study} --- so every
    study in Unit 1 so far has been one.}
  \end{center}
\end{frame}

% ── ROW 1b. THE TRAP ─────────────────────────────────────────────────────────
\begin{frame}[t, plain]
  \ceruleanheader{Problem 3 --- vote first, then decide}
  \sectionlabel{notes row 1 $\cdot$ the trap $\cdot$ ask \emph{who decided the groups}}
  \vspace{6pt}
  \begin{block}{``Assign half the members to swim at 6 a.m.\ and half at 7 p.m.\ for a month, then compare their sleep.''}
    Observational? --- It is not a survey, but that is not the question. \redink{Somebody
    assigned the groups}, so it is an experiment. Lesson 1.6.
  \end{block}
  \vspace{8pt}
  \begin{itemize}\setlength{\itemsep}{4pt}
    \item \textbf{2.} Count how many members walk in before 8 a.m.\ every day for a week.
          \redink{Observational} --- watching only, no questions asked.
    \item \textbf{4.} Mail the survey to \emph{every} member --- a census. \redink{Still
          observational.} A census asks everyone; it still assigns nobody.
  \end{itemize}
  \vspace{6pt}
  \begin{center}
    {\color{cerulean}\bfseries The test is never ``is it a survey.'' It is \emph{who put each
    person in their group?}}
  \end{center}
\end{frame}

% ── ROW 2. THE CYCLE — PS.DC.2d proper ───────────────────────────────────────
\begin{frame}[t, plain]
  \ceruleanheader{Planning one with the statistical cycle}
  \sectionlabel{notes row 2 $\cdot$ The Cycle $\cdot$ problems 5--8 $\cdot$ PS.DC.2d}
  \vspace{2pt}
  \begin{center}
    \small
    \renewcommand{\arraystretch}{1.35}
    \begin{tabular}{@{} l l @{}}
    \hline
    \textbf{Stage} & \textbf{What the pool study does} \\
    \hline
    1 Formulate & Do morning swimmers report more sleep than evening swimmers? \\
    2 Collect   & Frame: all 1{,}500 members. Random sample of 200; all 200 answered. \\
    3 Represent & A two-way table: swim time by sleep (7+ hrs or fewer). \\
    4 Analyze   & 70\% vs.\ 45\% --- a 25-point gap, morning higher. \\
    \hline
    \end{tabular}
  \end{center}
  \vspace{8pt}
  \begin{block}{Problem 8 --- what the plan never says}
    Nowhere does anybody tell a member \emph{when to swim}. \redink{That one absence is the
    whole definition}, written as a plan. Lesson 1.3's sampling plan and Lesson 1.4's bias
    check are both still standing.
  \end{block}
\end{frame}

% ── ROW 3. TWO VARIABLES ─────────────────────────────────────────────────────
\begin{frame}[t, plain]
  \ceruleanheader{Two groups that built themselves}
  \sectionlabel{notes row 3 $\cdot$ Two Variables $\cdot$ problems 9--10}
  \vspace{2pt}
  \begin{center}
    \small
    \renewcommand{\arraystretch}{1.3}
    \begin{tabular}{@{} l c c c @{}}
    \hline
     & \textbf{7+ hrs} & \textbf{Fewer} & \textbf{Total} \\
    \hline
    Morning swimmers & 56 & 24 & \textbf{80} \\
    Evening swimmers & 54 & 66 & \textbf{120} \\
    \hline
    Total & 110 & 90 & \textbf{200} \\
    \hline
    \end{tabular}
  \end{center}
  \vspace{6pt}
  \begin{itemize}\setlength{\itemsep}{4pt}
    \item \textbf{9.} $56/80 =$ \redink{70\%} \quad $54/120 =$ \redink{45\%} \quad a
          \redink{25-point} gap, morning higher. \emph{Two counts, two different wholes.}
    \item \textbf{10.} Pooled: $110/200 =$ \redink{55\%}, so about \redink{825} of all 1{,}500
          members --- \emph{an estimate}.
  \end{itemize}
  \vspace{6pt}
  \begin{block}{The arrow in the claim}
    Tail = the \textbf{explanatory} variable (swim time). Head = the \textbf{response}
    variable (7+ hours of sleep). They are \textbf{associated}: knowing the group changes what
    you expect. \redink{Nobody built these groups --- the members chose their own swim time.}
  \end{block}
\end{frame}

% ── ROW 4a. THE CRUX — three stories ─────────────────────────────────────────
\begin{frame}[t, plain]
  \ceruleanheader{Three stories fit this one table}
  \sectionlabel[redacc]{notes row 4 $\cdot$ Not Causation $\cdot$ problems 11--14 $\cdot$ the whole point of today}
  \vspace{4pt}
  \begin{columns}[t]
    \begin{column}{0.31\textwidth}
      \begin{block}{1 The claim}
        morning swim $\to$ more sleep
      \end{block}
    \end{column}
    \begin{column}{0.31\textwidth}
      \begin{block}{2 Backwards}
        more sleep $\to$ morning swim. \emph{Who can get up at 6 a.m.?} (problem 11)
      \end{block}
    \end{column}
    \begin{column}{0.31\textwidth}
      \begin{block}{3 Lurking}
        \textbf{retired?} moves both. \redink{60/80 = 75\%} of morning vs.\ \redink{18/120 =
        15\%} of evening (problem 12)
      \end{block}
    \end{column}
  \end{columns}
  \vspace{10pt}
  \begin{center}
    {\large The 25-point gap is real, and it is \redink{not bias} --- the sample was random and
    everyone answered.}\\[6pt]
    {\small\color{slate} A \textbf{lurking variable} sits \emph{outside} the study and affects
    \emph{both} variables in it. Retirement was never in the table.}
  \end{center}
\end{frame}

% ── ROW 4b. THE VOTE AGAIN ───────────────────────────────────────────────────
\begin{frame}[t, plain]
  \ceruleanheader{Problem 13 --- vote again: is the flyer safe?}
  \sectionlabel[redacc]{notes row 4 $\cdot$ the crux settles here}
  \vspace{4pt}
  \begin{columns}[t]
    \begin{column}{0.47\textwidth}
      \begin{block}{``Morning swimmers report more sleep.''}
        \redink{Supported.} It is what was seen --- both percents and the gap.
      \end{block}
    \end{column}
    \begin{column}{0.47\textwidth}
      \begin{block}{``Swimming in the morning makes you sleep longer.''}
        \redink{Not supported.} \emph{Makes} is causation, and three stories fit the table.
      \end{block}
    \end{column}
  \end{columns}
  \vspace{8pt}
  \begin{block}{Problem 14 --- ``If the evening swimmers switched to mornings, they would sleep more.''}
    \redink{Not supported either.} A prediction about \emph{switching groups} is a causal claim
    in different clothes --- nobody assigned anyone to switch.
  \end{block}
  \vspace{4pt}
  \begin{center}
    {\color{cerulean}\bfseries An observational study can show \emph{association}. It cannot
    show \emph{causation}, because nobody assigned the groups.}\\[4pt]
    {\small\color{slate} Banned verbs: \emph{causes, makes, raises, improves, helps, hurts.}}
  \end{center}
\end{frame}

% ── WE DO — a LIVE surface: task and data only, no answers ───────────────────
\begin{frame}[t, plain]
  \ceruleanheader{Guided Practice: the Riverbend Study Center}
  \sectionlabel[goldacc]{We do $\cdot$ problems 15--18 $\cdot$ you hold the pen}
  Riverbend has \textbf{900} students and a Study Center anyone may walk into. A counselor
  samples \textbf{150} records at random. Of the \textbf{60} who used it weekly, \textbf{39}
  passed every class; of the \textbf{90} who never did, \textbf{72} did.
  \vspace{4pt}
  \begin{itemize}\setlength{\itemsep}{4pt}
    \item[15.] Observational? Who decided the groups? Explanatory and response variables?
    \item[16.] Both pass rates and the gap in percentage points. Which group is higher?
    \item[17.] Write the finding as an \emph{association} sentence --- no causal verb. Does it
          support the memo?
    \item[18.] Last quarter, before the Center opened: 45 of the 60 users had already failed a
          class; 9 of the 90 non-users had. Percents? Which story is this?
  \end{itemize}
  \begin{block}{The principal's memo}
    \emph{``The Study Center is hurting grades. Shut it down.''} --- Which came first?
  \end{block}
\end{frame}

% ── GUIDED PRACTICE problem 18 — shown AFTER the class has worked it ─────────
\begin{frame}[t, plain]
  \ceruleanheader{The arrow was backwards the whole time}
  \sectionlabel[goldacc]{Guided Practice, problem 18 $\cdot$ shown after you have worked it}
  \vspace{4pt}
  \begin{center}
    \small
    \renewcommand{\arraystretch}{1.3}
    \begin{tabular}{@{} l c c @{}}
    \hline
    & \textbf{Passed every class} & \textbf{Had already failed a class \emph{last} quarter} \\
    \hline
    Used the Center (60)  & $39/60 = 65\%$ & $45/60 = 75\%$ \\
    Never used it (90)    & $72/90 = 80\%$ & $9/90 = 10\%$ \\
    \hline
    \end{tabular}
  \end{center}
  \vspace{8pt}
  \begin{block}{What the study supports}
    \emph{``At Riverbend, 65\% of Study Center users passed every class against 80\% of
    non-users --- using the Center is \textbf{associated} with a lower pass rate.''}
    \redink{The memo is not supported.} The low grades sent the students to the Center; the
    Center did not lower the grades. \textbf{It looks worse than it is.}
  \end{block}
\end{frame}

% ── DEBRIEF ──────────────────────────────────────────────────────────────────
\begin{frame}[t, plain]
  \ceruleanheader{Debrief: three stories, one table}
  \sectionlabel{10 minutes $\cdot$ pens down, papers up --- we say these aloud}
  \vspace{4pt}
  \begin{enumerate}\setlength{\itemsep}{4pt}
    \item The claim, the arrow backwards, retirement --- say the sentence. The study cannot
          tell them apart.
    \item Two problem-17 answers side by side. The difference is one verb.
    \item Problem 18: which came first, and does the Center look better or worse than it is?
  \end{enumerate}
  \vspace{6pt}
  \begin{block}{Cold check --- hands down, thirty seconds}
    Bayside Middle School. Of the \textbf{150} students in band, \textbf{90} earned an A or B
    in math; of the \textbf{600} not in band, \textbf{240} did. The principal concludes band
    \emph{raises} math grades and wants it required. \textbf{Is the principal right?}
    \quad {\small\color{slate}(60\% vs.\ 40\%, a 20-point gap --- an association. Nobody was
    placed in band; students signed up.)}
  \end{block}
\end{frame}

% ── YOU DO — the homework is the individual practice, started in class ───────
\begin{frame}[t, plain]
  \ceruleanheader{You do: start the homework}
  \sectionlabel[goldacc]{Last 10 minutes $\cdot$ on your own $\cdot$ finish it at home}
  \begin{itemize}\setlength{\itemsep}{4pt}
    \item \textbf{Millbrook Public Library} --- 1{,}200 card holders, an unbiased phone sample
          of 240, everyone reached. Does the Summer Reading Program work?
    \item Name the study and the two variables; both rates with \emph{two different}
          denominators; pool the sample and scale it to 1{,}200.
    \item Two sentences on the same data --- decide which one the study supports, and why.
    \item The previous summer's records, and the direction of the error.
    \item \textbf{No bias at all --- and still no causation.} Explain why.
  \end{itemize}
  \begin{block}{Silent and alone --- I am circulating. Packet pages, scored out of 12, due at the start of the first class after two study halls.}
    Start with \textbf{1, 2, and 4}. Item 6 is the one I am reading over your shoulder --- if
    your answer is ``the study was biased,'' flag me before you leave.
  \end{block}
\end{frame}

% ── CLOSE ────────────────────────────────────────────────────────────────────
{
\setbeamercolor{background canvas}{bg=cerulean}
\begin{frame}[plain]
  \vspace{1.5cm}\hspace{0.8cm}%
  \begin{minipage}{0.86\textwidth}
    {\color{goldacc}\bfseries\small BEFORE YOU GO}\\[0.7cm]
    {\color{white}\bfseries\Large In an observational study you watch --- you do not assign. A
    random sample, a 100\% response rate, and perfect arithmetic still do not buy you the word
    \emph{causes}.}\\[0.9cm]
    {\color{frostmid}\small Next: Lesson 1.6 --- \emph{Experimental Design}. Today nobody
    assigned the groups, so three stories fit one table. Next time we assign them ourselves.}
  \end{minipage}
\end{frame}
}

\end{document}
